\documentclass[11pt,a4paper]{article}
\usepackage[utf8]{inputenc}
\usepackage[T1]{fontenc}
\usepackage[catalan]{babel}
\usepackage{lmodern}
\usepackage[sfdefault,lining]{FiraSans}
\usepackage[a4paper,top=2.3cm,bottom=2.2cm,left=1.8cm,right=1.8cm,headheight=1.3cm,headsep=0.4cm,footskip=1cm]{geometry}
\usepackage{amsmath,amssymb}
\usepackage{icomma}
\usepackage{graphicx}
\usepackage[table]{xcolor}
\usepackage{tikz}
\usepackage{fontawesome5}
\usepackage{ccicons}
\usepackage{fancyhdr}
\usepackage{lastpage}
\usepackage{tasks}
\usepackage{needspace}
\usepackage[most]{tcolorbox}
\usepackage{array}
\usepackage[hidelinks,pdftitle={Fitxa d'exercicis de la unitat 2: Decimals i fraccions. 2n d'ESO},pdfsubject={Matemàtiques 2n ESO},pdfkeywords={nombres decimals, decimals exactes i periòdics, operacions amb decimals, fraccions equivalents, fracció generatriu, operacions amb fraccions, potències de fraccions, operacions combinades, problemes amb fraccions}]{hyperref}

\definecolor{navy}{HTML}{1B2748}
\definecolor{accent}{HTML}{E8603A}
\definecolor{teal}{HTML}{1F9E90}
\definecolor{soft}{HTML}{EEF1F8}
\definecolor{muted}{HTML}{5D6685}

\setlength{\parindent}{0pt}
\setlength{\parskip}{3pt}
\settasks{label=\alph*), label-format=\bfseries\color{navy}, label-width=1.3em, item-indent=1.8em, column-sep=0.6em, after-item-skip=5pt, before-skip=3pt, after-skip=4pt}

% ---------- capçalera i peu ----------
\newcommand{\llicencia}{\href{https://creativecommons.org/licenses/by-nc-sa/4.0/deed.ca}{\ccbyncsa}\ \ Llicència Creative Commons Reconeixement-NoComercial-CompartirIgual 4.0 Internacional (CC BY-NC-SA 4.0)}
\pagestyle{fancy}
\fancyhf{}
\fancyhead[L]{\raisebox{-0.35\height}{\includegraphics[height=0.95cm]{logo-gva.png}}}
\fancyhead[C]{\small\color{muted}Fitxa d'exercicis · Unitat 2: Decimals i fraccions · 2n d'ESO}
\fancyhead[R]{\raisebox{-0.35\height}{\includegraphics[height=0.95cm]{logo-vermellar.png}}}
\fancyfoot[L]{\scriptsize\color{muted}\llicencia}
\fancyfoot[R]{\small\color{muted}\thepage}
\renewcommand{\headrulewidth}{0.4pt}
\renewcommand{\headrule}{\color{muted!40}\hrule height \headrulewidth}
\fancypagestyle{primera}{\fancyhf{}\renewcommand{\headrulewidth}{0pt}
  \fancyfoot[L]{\scriptsize\color{muted}\llicencia}
  \fancyfoot[R]{\small\color{muted}\thepage}}

% ---------- nivells i exercicis ----------
\newcommand{\estrelles}[1]{{\color{accent}\ifcase#1\or\faStar\faStar[regular]\faStar[regular]\or\faStar\faStar\faStar[regular]\or\faStar\faStar\faStar\fi}}
\newcounter{exer}
\newcommand{\ex}[1]{\par\needspace{4\baselineskip}\medskip\stepcounter{exer}\noindent{\color{navy}\bfseries\theexer.}\ \ {\small\estrelles{#1}}\ \ }

\newcommand{\seccio}[2]{\par\needspace{10\baselineskip}\bigskip
  \begin{tcolorbox}[colback=navy,colframe=navy,coltext=white,arc=2mm,boxrule=0pt,left=3mm,right=3mm,top=1.2mm,bottom=1.2mm]
  {\large\bfseries #1\quad #2}
  \end{tcolorbox}\vspace{-2pt}}

\newcolumntype{C}[1]{>{\centering\arraybackslash}p{#1}}
% quadre de recordatori (només a la fitxa d'exercicis)
\newcommand{\recorda}[1]{\par\smallskip\begin{tcolorbox}[colback=accent!7,colframe=accent!55,arc=2mm,boxrule=0.5pt,left=2.5mm,right=2.5mm,top=1mm,bottom=1mm]\small{\color{accent}\faLightbulb[regular]}\ \textbf{Recorda:}\ #1\end{tcolorbox}}
\newcommand{\res}[1]{#1}\newcommand{\rest}[1]{#1}

\def\fmneed{12}
% =====================================================================
% Capçalera FentMates de les fitxes: un color per a cada curs.
% Els colors són els de la pàgina inicial (LEVELS a assets/js/materials.js);
% el logo ve de _eines/logo/pdf/ (python3 _eines/logo/genera.py).
%
% Ús, al final del preàmbul de la unitat:
%   % =====================================================================
% Capçalera FentMates de les fitxes: un color per a cada curs.
% Els colors són els de la pàgina inicial (LEVELS a assets/js/materials.js);
% el logo ve de _eines/logo/pdf/ (python3 _eines/logo/genera.py).
%
% Ús, al final del preàmbul de la unitat:
%   % =====================================================================
% Capçalera FentMates de les fitxes: un color per a cada curs.
% Els colors són els de la pàgina inicial (LEVELS a assets/js/materials.js);
% el logo ve de _eines/logo/pdf/ (python3 _eines/logo/genera.py).
%
% Ús, al final del preàmbul de la unitat:
%   \input{../comu/fentmates.tex}
%   \fmcurs{1eso}
% i, al començament del document:
%   \capcalera{Matemàtiques}{1r ESO}{Fitxa d'exercicis · Unitat 2}{Divisibilitat}
% Opcional abans de \input: \def\fmneed{14} (espai mínim abans de cada secció).
% =====================================================================
\makeatletter
\def\fm@logodir{../../logo/pdf}
\providecommand{\fmneed}{10}
\def\fm@cursnom{}

% clau del curs, color i color del text damunt del color
\newcommand{\fm@curs}[3]{\@namedef{fm@color@#1}{#2}\@namedef{fm@text@#1}{#3}}
\fm@curs{1eso}{E6B84A}{navy}
\fm@curs{2eso}{3B82F6}{white}
\fm@curs{3eso}{FF7A50}{white}
\fm@curs{4eso-a}{A78BFA}{white}
\fm@curs{4eso-b}{6D28D9}{white}
\fm@curs{1batx-mi}{9B1C1C}{white}
\fm@curs{1batx-ccss}{EF6B6B}{white}
\fm@curs{2batx-mii}{1E6B47}{white}
\fm@curs{2batx-ccss}{2BB3A3}{white}

\newcommand{\fmcurs}[1]{%
  \@ifundefined{fm@color@#1}{\PackageError{fentmates}{Curs desconegut: #1}{}}{}%
  \edef\fm@tmp{\noexpand\definecolor{curs}{HTML}{\@nameuse{fm@color@#1}}%
               \noexpand\colorlet{curstext}{\@nameuse{fm@text@#1}}}\fm@tmp
  \colorlet{cursfosc}{curs!72!black}%
  \def\fm@logo{\fm@logodir/vertical-#1.pdf}%
  \def\fm@marca{\fm@logodir/marca-#1.pdf}}

% ---------- capçalera de la primera pàgina ----------
% #1 matèria (petit), #2 curs, #3 tipus de document i unitat, #4 títol
\newcommand{\capcalera}[4]{%
  \gdef\fm@cursnom{#2}%
  \par\noindent
  \begin{tcolorbox}[colback=curs,colframe=curs,boxrule=0pt,arc=2.2mm,width=3.5cm,height=2.15cm,
      valign=center,left=3mm,right=2mm,top=1mm,bottom=1mm,nobeforeafter,box align=center]
    \color{curstext}{\footnotesize #1}\par\vspace{1pt}{\fontsize{19}{21}\selectfont\bfseries #2\par}%
  \end{tcolorbox}\hspace{5mm}%
  \begin{minipage}[c]{\dimexpr\linewidth-3.5cm-5mm-2.9cm\relax}
    {\color{cursfosc}\small\bfseries\MakeUppercase{#3}}\par\vspace{2pt}
    {\color{navy}\fontsize{20}{23}\selectfont\bfseries #4\par}
  \end{minipage}\hfill
  \begin{minipage}[c]{2.6cm}\raggedleft\includegraphics[width=2.6cm]{\fm@logo}\end{minipage}\par
  \vspace{2pt}}

% ---------- barres de secció del color del curs ----------
\renewcommand{\seccio}[2]{\par\needspace{\fmneed\baselineskip}\bigskip
  \begin{tcolorbox}[colback=curs,colframe=curs,coltext=curstext,arc=2mm,boxrule=0pt,left=3mm,right=3mm,top=1.2mm,bottom=1.2mm]
  {\large\bfseries #1\quad #2}
  \end{tcolorbox}\vspace{-2pt}}

% ---------- estrelles de dificultat del color del curs (a Batxillerat no se'n posen) ----------
\@ifundefined{estrelles}{}{%
  \renewcommand{\estrelles}[1]{{\color{curs}\ifcase#1\or\faStar\faStar[regular]\faStar[regular]\or\faStar\faStar\faStar[regular]\or\faStar\faStar\faStar\fi}}}

% ---------- capçalera i peu de les pàgines ----------
% primera pàgina: logos institucionals al peu
\fancypagestyle{primera}{\fancyhf{}\renewcommand{\headrulewidth}{0pt}
  \fancyfoot[L]{\raisebox{-0.32\height}{\includegraphics[height=0.78cm]{logo-gva.png}}\hspace{5mm}%
    \raisebox{-0.32\height}{\includegraphics[height=0.78cm]{logo-vermellar.png}}}
  \fancyfoot[R]{\scriptsize\color{muted}\href{https://creativecommons.org/licenses/by-nc-sa/4.0/deed.ca}{\ccbyncsa}\ \ CC BY-NC-SA 4.0 · fentmates.org\qquad\small\thepage}}
% la resta: el símbol a l'esquerra i el curs a la dreta, amb el títol al centre com abans
\fancyhead[L]{\raisebox{-0.3\height}{\includegraphics[height=0.62cm]{\fm@marca}}}
\fancyhead[R]{\small\bfseries\color{cursfosc}\fm@cursnom}
\renewcommand{\headrule}{\color{curs!45}\hrule height \headrulewidth}
\makeatother

%   \fmcurs{1eso}
% i, al començament del document:
%   \capcalera{Matemàtiques}{1r ESO}{Fitxa d'exercicis · Unitat 2}{Divisibilitat}
% Opcional abans de \input: \def\fmneed{14} (espai mínim abans de cada secció).
% =====================================================================
\makeatletter
\def\fm@logodir{../../logo/pdf}
\providecommand{\fmneed}{10}
\def\fm@cursnom{}

% clau del curs, color i color del text damunt del color
\newcommand{\fm@curs}[3]{\@namedef{fm@color@#1}{#2}\@namedef{fm@text@#1}{#3}}
\fm@curs{1eso}{E6B84A}{navy}
\fm@curs{2eso}{3B82F6}{white}
\fm@curs{3eso}{FF7A50}{white}
\fm@curs{4eso-a}{A78BFA}{white}
\fm@curs{4eso-b}{6D28D9}{white}
\fm@curs{1batx-mi}{9B1C1C}{white}
\fm@curs{1batx-ccss}{EF6B6B}{white}
\fm@curs{2batx-mii}{1E6B47}{white}
\fm@curs{2batx-ccss}{2BB3A3}{white}

\newcommand{\fmcurs}[1]{%
  \@ifundefined{fm@color@#1}{\PackageError{fentmates}{Curs desconegut: #1}{}}{}%
  \edef\fm@tmp{\noexpand\definecolor{curs}{HTML}{\@nameuse{fm@color@#1}}%
               \noexpand\colorlet{curstext}{\@nameuse{fm@text@#1}}}\fm@tmp
  \colorlet{cursfosc}{curs!72!black}%
  \def\fm@logo{\fm@logodir/vertical-#1.pdf}%
  \def\fm@marca{\fm@logodir/marca-#1.pdf}}

% ---------- capçalera de la primera pàgina ----------
% #1 matèria (petit), #2 curs, #3 tipus de document i unitat, #4 títol
\newcommand{\capcalera}[4]{%
  \gdef\fm@cursnom{#2}%
  \par\noindent
  \begin{tcolorbox}[colback=curs,colframe=curs,boxrule=0pt,arc=2.2mm,width=3.5cm,height=2.15cm,
      valign=center,left=3mm,right=2mm,top=1mm,bottom=1mm,nobeforeafter,box align=center]
    \color{curstext}{\footnotesize #1}\par\vspace{1pt}{\fontsize{19}{21}\selectfont\bfseries #2\par}%
  \end{tcolorbox}\hspace{5mm}%
  \begin{minipage}[c]{\dimexpr\linewidth-3.5cm-5mm-2.9cm\relax}
    {\color{cursfosc}\small\bfseries\MakeUppercase{#3}}\par\vspace{2pt}
    {\color{navy}\fontsize{20}{23}\selectfont\bfseries #4\par}
  \end{minipage}\hfill
  \begin{minipage}[c]{2.6cm}\raggedleft\includegraphics[width=2.6cm]{\fm@logo}\end{minipage}\par
  \vspace{2pt}}

% ---------- barres de secció del color del curs ----------
\renewcommand{\seccio}[2]{\par\needspace{\fmneed\baselineskip}\bigskip
  \begin{tcolorbox}[colback=curs,colframe=curs,coltext=curstext,arc=2mm,boxrule=0pt,left=3mm,right=3mm,top=1.2mm,bottom=1.2mm]
  {\large\bfseries #1\quad #2}
  \end{tcolorbox}\vspace{-2pt}}

% ---------- estrelles de dificultat del color del curs (a Batxillerat no se'n posen) ----------
\@ifundefined{estrelles}{}{%
  \renewcommand{\estrelles}[1]{{\color{curs}\ifcase#1\or\faStar\faStar[regular]\faStar[regular]\or\faStar\faStar\faStar[regular]\or\faStar\faStar\faStar\fi}}}

% ---------- capçalera i peu de les pàgines ----------
% primera pàgina: logos institucionals al peu
\fancypagestyle{primera}{\fancyhf{}\renewcommand{\headrulewidth}{0pt}
  \fancyfoot[L]{\raisebox{-0.32\height}{\includegraphics[height=0.78cm]{logo-gva.png}}\hspace{5mm}%
    \raisebox{-0.32\height}{\includegraphics[height=0.78cm]{logo-vermellar.png}}}
  \fancyfoot[R]{\scriptsize\color{muted}\href{https://creativecommons.org/licenses/by-nc-sa/4.0/deed.ca}{\ccbyncsa}\ \ CC BY-NC-SA 4.0 · fentmates.org\qquad\small\thepage}}
% la resta: el símbol a l'esquerra i el curs a la dreta, amb el títol al centre com abans
\fancyhead[L]{\raisebox{-0.3\height}{\includegraphics[height=0.62cm]{\fm@marca}}}
\fancyhead[R]{\small\bfseries\color{cursfosc}\fm@cursnom}
\renewcommand{\headrule}{\color{curs!45}\hrule height \headrulewidth}
\makeatother

%   \fmcurs{1eso}
% i, al començament del document:
%   \capcalera{Matemàtiques}{1r ESO}{Fitxa d'exercicis · Unitat 2}{Divisibilitat}
% Opcional abans de \input: \def\fmneed{14} (espai mínim abans de cada secció).
% =====================================================================
\makeatletter
\def\fm@logodir{../../logo/pdf}
\providecommand{\fmneed}{10}
\def\fm@cursnom{}

% clau del curs, color i color del text damunt del color
\newcommand{\fm@curs}[3]{\@namedef{fm@color@#1}{#2}\@namedef{fm@text@#1}{#3}}
\fm@curs{1eso}{E6B84A}{navy}
\fm@curs{2eso}{3B82F6}{white}
\fm@curs{3eso}{FF7A50}{white}
\fm@curs{4eso-a}{A78BFA}{white}
\fm@curs{4eso-b}{6D28D9}{white}
\fm@curs{1batx-mi}{9B1C1C}{white}
\fm@curs{1batx-ccss}{EF6B6B}{white}
\fm@curs{2batx-mii}{1E6B47}{white}
\fm@curs{2batx-ccss}{2BB3A3}{white}

\newcommand{\fmcurs}[1]{%
  \@ifundefined{fm@color@#1}{\PackageError{fentmates}{Curs desconegut: #1}{}}{}%
  \edef\fm@tmp{\noexpand\definecolor{curs}{HTML}{\@nameuse{fm@color@#1}}%
               \noexpand\colorlet{curstext}{\@nameuse{fm@text@#1}}}\fm@tmp
  \colorlet{cursfosc}{curs!72!black}%
  \def\fm@logo{\fm@logodir/vertical-#1.pdf}%
  \def\fm@marca{\fm@logodir/marca-#1.pdf}}

% ---------- capçalera de la primera pàgina ----------
% #1 matèria (petit), #2 curs, #3 tipus de document i unitat, #4 títol
\newcommand{\capcalera}[4]{%
  \gdef\fm@cursnom{#2}%
  \par\noindent
  \begin{tcolorbox}[colback=curs,colframe=curs,boxrule=0pt,arc=2.2mm,width=3.5cm,height=2.15cm,
      valign=center,left=3mm,right=2mm,top=1mm,bottom=1mm,nobeforeafter,box align=center]
    \color{curstext}{\footnotesize #1}\par\vspace{1pt}{\fontsize{19}{21}\selectfont\bfseries #2\par}%
  \end{tcolorbox}\hspace{5mm}%
  \begin{minipage}[c]{\dimexpr\linewidth-3.5cm-5mm-2.9cm\relax}
    {\color{cursfosc}\small\bfseries\MakeUppercase{#3}}\par\vspace{2pt}
    {\color{navy}\fontsize{20}{23}\selectfont\bfseries #4\par}
  \end{minipage}\hfill
  \begin{minipage}[c]{2.6cm}\raggedleft\includegraphics[width=2.6cm]{\fm@logo}\end{minipage}\par
  \vspace{2pt}}

% ---------- barres de secció del color del curs ----------
\renewcommand{\seccio}[2]{\par\needspace{\fmneed\baselineskip}\bigskip
  \begin{tcolorbox}[colback=curs,colframe=curs,coltext=curstext,arc=2mm,boxrule=0pt,left=3mm,right=3mm,top=1.2mm,bottom=1.2mm]
  {\large\bfseries #1\quad #2}
  \end{tcolorbox}\vspace{-2pt}}

% ---------- estrelles de dificultat del color del curs (a Batxillerat no se'n posen) ----------
\@ifundefined{estrelles}{}{%
  \renewcommand{\estrelles}[1]{{\color{curs}\ifcase#1\or\faStar\faStar[regular]\faStar[regular]\or\faStar\faStar\faStar[regular]\or\faStar\faStar\faStar\fi}}}

% ---------- capçalera i peu de les pàgines ----------
% primera pàgina: logos institucionals al peu
\fancypagestyle{primera}{\fancyhf{}\renewcommand{\headrulewidth}{0pt}
  \fancyfoot[L]{\raisebox{-0.32\height}{\includegraphics[height=0.78cm]{logo-gva.png}}\hspace{5mm}%
    \raisebox{-0.32\height}{\includegraphics[height=0.78cm]{logo-vermellar.png}}}
  \fancyfoot[R]{\scriptsize\color{muted}\href{https://creativecommons.org/licenses/by-nc-sa/4.0/deed.ca}{\ccbyncsa}\ \ CC BY-NC-SA 4.0 · fentmates.org\qquad\small\thepage}}
% la resta: el símbol a l'esquerra i el curs a la dreta, amb el títol al centre com abans
\fancyhead[L]{\raisebox{-0.3\height}{\includegraphics[height=0.62cm]{\fm@marca}}}
\fancyhead[R]{\small\bfseries\color{cursfosc}\fm@cursnom}
\renewcommand{\headrule}{\color{curs!45}\hrule height \headrulewidth}
\makeatother

\fmcurs{2eso}

\begin{document}
\thispagestyle{primera}

\capcalera{Matemàtiques}{2n ESO}{Fitxa d'exercicis · Unitat 2}{Decimals i fraccions}
\vspace{10pt}
Nom i cognoms: \rule{8.4cm}{0.4pt}\hfill Grup: \rule{1.6cm}{0.4pt}\hfill Data: \rule{2.2cm}{0.4pt}
\vspace{8pt}

\begin{tcolorbox}[colback=soft,colframe=soft,arc=2mm,left=3mm,right=3mm,top=1.5mm,bottom=1.5mm]
\small Exercicis graduats per nivells de dificultat:\quad \estrelles{1}\ bàsic \qquad \estrelles{2}\ mitjà \qquad \estrelles{3}\ avançat.\\[2pt] Resol-los al quadern mostrant tots els passos i sense calculadora. Dona les fraccions en forma irreductible. Les seccions segueixen l'ordre dels apunts de la unitat (2.1--2.11).
\end{tcolorbox}

% =====================================================
\seccio{1}{Nombres decimals}

\ex{1} Escriu com es llegeixen aquests nombres:\begin{tasks}(3)
  \task $7{,}25$
  \task $0{,}408$
  \task $12{,}06$
  \task $3{,}5$
  \task $100{,}001$
  \task $0{,}09$
\end{tasks}\recorda{llegim la part entera, la paraula \emph{unitats} i la part decimal amb el nom de l'\textbf{última} xifra: $3{,}045$ són 3 unitats i 45 mil·lèsimes.}

\ex{1} Quant val la xifra indicada en cada nombre?\begin{tasks}(3)
  \task El 7 en $34{,}571$
  \task El 3 en $2{,}738$
  \task El 7 en $607{,}15$
  \task El 8 en $0{,}0809$
  \task El 4 en $15{,}4$
  \task El 6 en $9{,}206$
\end{tasks}

\ex{1} Escriu $<$, $>$ o $=$:\begin{tasks}(3)
  \task $4{,}7\ \square\ 4{,}68$
  \task $0{,}09\ \square\ 0{,}1$
  \task $5{,}307\ \square\ 5{,}37$
  \task $12{,}5\ \square\ 12{,}5$
  \task $3{,}099\ \square\ 3{,}1$
  \task $0{,}45\ \square\ 0{,}405$
\end{tasks}

\ex{2} Ordena de menor a major: $2{,}5;\ 2{,}05;\ 2{,}55;\ 2{,}505;\ 2{,}15;\ 2{,}501$

\ex{2} Completa la taula:\par\smallskip{\centering\renewcommand{\arraystretch}{1.35}\begin{tabular}{|c|C{2.4cm}|C{2.4cm}|C{2.4cm}|}\hline Nombre & Truncat a les centèsimes & Arrodonit a les dècimes & Arrodonit a les centèsimes\\\hline $3{,}4682$ & & & \\\hline $15{,}8962$ & & & \\\hline $0{,}7049$ & & & \\\hline $29{,}996$ & & & \\\hline\end{tabular}\par}\recorda{per a \textbf{arrodonir}, mirem la primera xifra que eliminem: si és 0, 1, 2, 3 o 4, l'última xifra es queda igual; si és 5, 6, 7, 8 o 9, li sumem 1.}

\ex{3} Escriu tres nombres compresos entre $4{,}7$ i $4{,}71$. Quin és el nombre més xicotet amb dues xifres decimals que, arrodonit a les unitats, dona 8? I el més gran?


% =====================================================
\seccio{2}{Tipus de nombres decimals}

\ex{1} Classifica els nombres següents i, si són periòdics, escriu-los amb la ratlla del període:\begin{tasks}(2)
  \task $4{,}252525\ldots$
  \task $0{,}125$
  \task $7{,}31666\ldots$
  \task $1{,}010010001\ldots$
  \task $0{,}777\ldots$
  \task $12{,}408408\ldots$
  \task $3{,}1252525\ldots$
  \task $\pi=3{,}14159\ldots$
\end{tasks}\recorda{periòdic \textbf{pur}: el període comença just després de la coma; \textbf{mixt}: abans del període hi ha xifres que no es repeteixen (l'anteperíode).}

\ex{1} Indica la part entera, l'anteperíode i el període de cada nombre:\begin{tasks}(4)
  \task $5{,}\overline{6}$
  \task $0{,}0\overline{9}$
  \task $2{,}13\overline{4}$
  \task $0{,}\overline{142857}$
\end{tasks}

\ex{2} Ordena de menor a major: $0{,}\overline{3};\ \ 0{,}3;\ \ 0{,}3\overline{4};\ \ 0{,}\overline{34};\ \ 0{,}33$

\ex{3} Escriu un nombre decimal periòdic pur, un de periòdic mixt, un d'exacte i un de no exacte i no periòdic que estiguen tots entre $2{,}1$ i $2{,}2$.


% =====================================================
\seccio{3}{Sumar, restar i multiplicar decimals}

\ex{1} Calcula:\begin{tasks}(3)
  \task $23{,}45+7{,}8$
  \task $40-12{,}35$
  \task $0{,}9+0{,}35$
  \task $15-6{,}47$
  \task $3{,}06+12{,}9+0{,}4$
  \task $100-0{,}01$
\end{tasks}\recorda{per a sumar i restar, col·loquem les \textbf{comes alineades} i completem amb zeros.}

\ex{1} Calcula:\begin{tasks}(3)
  \task $4{,}6\cdot 3$
  \task $2{,}35\cdot 0{,}4$
  \task $0{,}6\cdot 0{,}05$
  \task $12{,}5\cdot 1{,}2$
  \task $3{,}25\cdot 1{,}4$
  \task $0{,}07\cdot 0{,}3$
\end{tasks}\recorda{multipliquem com si no hi haguera comes i separem tants decimals com \textbf{sumen} els dos factors.}

\ex{1} Calcula mentalment:\begin{tasks}(4)
  \task $3{,}7\cdot 100$
  \task $0{,}052\cdot 1000$
  \task $45{,}2:1000$
  \task $8:100$
  \task $6{,}15\cdot 10$
  \task $0{,}3:10$
  \task $250\cdot 0{,}1$
  \task $1{,}4\cdot 0{,}01$
\end{tasks}

\ex{2} Compre $2{,}5$ kg de pomes a $1{,}86$ €/kg i $0{,}75$ kg de formatge a $14{,}20$ €/kg. Pague amb un bitllet de 20 €. Quant em tornen?

\ex{2} Un full de paper fa $0{,}012$ cm de gruix. Quant fa un paquet de 500 fulls? I una pila de 1000 paquets, en metres?

\ex{3} Troba el nombre que falta: \ a) $3{,}75+\square=10$ \quad b) $\square-2{,}08=5{,}3$ \quad c) $\square\cdot100=4{,}5$ \quad d) $0{,}4\cdot\square=0{,}12$


% =====================================================
\seccio{4}{Dividir amb decimals}

\ex{1} Calcula:\begin{tasks}(3)
  \task $8{,}4:4$
  \task $15:8$
  \task $9{,}75:3$
  \task $4{,}5:0{,}9$
  \task $2{,}73:0{,}7$
  \task $6:0{,}12$
\end{tasks}\recorda{si el divisor té decimals, multipliquem dividend i divisor per 10, 100, 1000… fins que el divisor siga un nombre natural: $4{,}5:0{,}9=45:9$.}

\ex{2} Calcula el quocient amb les xifres decimals indicades (arrodoneix l'última):\begin{tasks}(4)
  \task $10:3$ \ (2)
  \task $7:6$ \ (2)
  \task $2:7$ \ (3)
  \task $5{,}3:0{,}9$ \ (2)
\end{tasks}

\ex{2} Calcula respectant la jerarquia de les operacions:\begin{tasks}[after-item-skip=6pt](2)
  \task $5{,}2-1{,}5\cdot 2{,}4+0{,}6:0{,}2$
  \task $\left(3{,}5-1{,}25\right)\cdot 0{,}4+2{,}7:0{,}09$
  \task $1{,}2\cdot \left(4-2{,}5\right)+0{,}8:4$
  \task $10-{\left(0{,}5+1{,}5\right)}^{2}\cdot 2$
\end{tasks}

\ex{2} Un tros de formatge de $0{,}85$ kg costa $11{,}39$ €. Quin és el preu del quilo? Quant costaria un tros de $1{,}2$ kg?

\ex{3} Un cotxe gasta $31{,}5$ L de gasolina en 450 km. Quants litres gasta cada 100 km? Si el litre costa $1{,}58$ €, quant costa un viatge de 260 km?

\ex{3} Amb una ampolla de $1{,}5$ L, quants gots de $0{,}25$ L podem omplir? I quants de $0{,}2$ L? Quant sobra en cada cas?


% =====================================================
\seccio{5}{Fraccions}

\ex{1} Calcula:\begin{tasks}(4)
  \task $\dfrac{2}{7}$ de $63$
  \task $\dfrac{3}{4}$ de $80$
  \task $\dfrac{5}{6}$ de $108$
  \task $\dfrac{3}{8}$ de $240$
\end{tasks}

\ex{1} Simplifica fins a obtindre la fracció irreductible:\begin{tasks}(6)
  \task $\dfrac{84}{126}$
  \task $\dfrac{105}{180}$
  \task $\dfrac{36}{48}$
  \task $\dfrac{75}{100}$
  \task $\dfrac{240}{360}$
  \task $\dfrac{98}{140}$
\end{tasks}\recorda{dividim numerador i denominador pel seu màxim comú divisor (MCD).}

\ex{2} Comprova si són equivalents:\begin{tasks}(4)
  \task $\dfrac{14}{21}\ \text{i}\ \dfrac{22}{33}$
  \task $\dfrac{3}{4}\ \text{i}\ \dfrac{12}{15}$
  \task $\dfrac{-6}{8}\ \text{i}\ \dfrac{9}{-12}$
  \task $\dfrac{5}{7}\ \text{i}\ \dfrac{15}{20}$
\end{tasks}

\ex{2} Troba $x$ perquè les fraccions siguen equivalents:\begin{tasks}(4)
  \task $\dfrac{x}{12}=\dfrac{15}{36}$
  \task $\dfrac{8}{x}=\dfrac{-2}{5}$
  \task $\dfrac{6}{9}=\dfrac{x}{15}$
  \task $\dfrac{3}{5}=\dfrac{27}{x}$
\end{tasks}

\ex{2} Ordena de menor a major: $\dfrac{5}{6};\ \ -\dfrac{3}{4};\ \ \dfrac{7}{9};\ \ -\dfrac{2}{3};\ \ \dfrac{3}{4}$

\ex{3} Escriu tres fraccions compreses entre $\dfrac{1}{3}$ i $\dfrac{1}{2}$. Quina fracció de denominador 20 està just a la meitat entre $\dfrac{2}{5}$ i $\dfrac{1}{2}$?


% =====================================================
\seccio{6}{Fraccions i decimals}

\ex{1} Sense fer la divisió, digues si cada fracció dona un decimal exacte, periòdic pur o periòdic mixt. Després, comprova-ho dividint:\begin{tasks}(4)
  \task $\dfrac{9}{20}$
  \task $\dfrac{5}{9}$
  \task $\dfrac{7}{30}$
  \task $\dfrac{21}{28}$
  \task $\dfrac{4}{11}$
  \task $\dfrac{7}{12}$
  \task $\dfrac{13}{40}$
  \task $\dfrac{1}{6}$
\end{tasks}\recorda{mira el denominador de la fracció \textbf{irreductible}: si només té els factors 2 i 5, és exacte; si no en té cap, periòdic pur; si en té i també d'altres, periòdic mixt.}

\ex{1} Escriu com a fracció irreductible:\begin{tasks}(6)
  \task $0{,}8$
  \task $0{,}45$
  \task $1{,}2$
  \task $0{,}004$
  \task $2{,}35$
  \task $0{,}125$
\end{tasks}

\ex{2} Calcula la fracció generatriu i simplifica-la:\begin{tasks}(4)
  \task $0{,}\overline{6}$
  \task $3{,}\overline{12}$
  \task $0{,}8\overline{3}$
  \task $2{,}1\overline{45}$
  \task $5{,}\overline{4}$
  \task $0{,}\overline{27}$
  \task $0{,}1\overline{6}$
  \task $1{,}2\overline{3}$
\end{tasks}\recorda{numerador: el nombre sense coma fins al final del període menys la part que no es repeteix; denominador: tants 9 com xifres té el període i tants 0 com xifres té l'anteperíode. \ $1{,}2\overline{3}=\dfrac{123-12}{90}$}

\ex{2} Comprova, dividint, que $\dfrac{37}{30}=1{,}2\overline{3}$. Després, demostra-ho multiplicant $x=1{,}2\overline{3}$ per 100 i per 10 i restant.

\ex{2} Ordena de menor a major passant-ho tot a decimal: \ $0{,}7;\ \ \dfrac{2}{3};\ \ 0{,}67;\ \ \dfrac{3}{4};\ \ 0{,}6\overline{8}$

\ex{3} Calcula passant a fracció: \ a) $0{,}\overline{3}+0{,}\overline{6}$ \quad b) $0{,}\overline{9}$ \quad c) $1{,}\overline{3}\cdot0{,}75$ \quad d) $0{,}1\overline{6}:0{,}\overline{5}$. Què observes en a) i b)?


% =====================================================
\seccio{7}{Suma i resta de fraccions}

\ex{1} Calcula i simplifica:\begin{tasks}[after-item-skip=6pt](3)
  \task $\dfrac{4}{7}+\dfrac{8}{7}$
  \task $\dfrac{3}{10}+\dfrac{2}{5}$
  \task $\dfrac{5}{6}-\dfrac{1}{4}$
  \task $\dfrac{7}{12}+\dfrac{5}{18}$
  \task $3-\dfrac{5}{4}$
  \task $1+\dfrac{2}{7}$
\end{tasks}\recorda{si els denominadors són diferents, reduïm a comú denominador amb el mcm: $\dfrac{5}{6}-\dfrac{1}{4}=\dfrac{10}{12}-\dfrac{3}{12}$.}

\ex{2} Calcula i simplifica:\begin{tasks}[after-item-skip=6pt](3)
  \task $\dfrac{1}{2}-\dfrac{2}{3}+\dfrac{1}{6}$
  \task $-\dfrac{3}{4}-\dfrac{1}{6}+2$
  \task $\dfrac{5}{6}-\dfrac{3}{4}+\dfrac{1}{3}$
  \task $\dfrac{2}{5}-\dfrac{7}{10}-1$
  \task $-\dfrac{1}{2}-\dfrac{1}{3}-\dfrac{1}{4}$
  \task $\dfrac{3}{8}-\dfrac{5}{12}+\dfrac{1}{6}$
\end{tasks}

\ex{2} Calcula (primer els parèntesis):\begin{tasks}[after-item-skip=8pt](2)
  \task $\dfrac{2}{5}-\left(\dfrac{1}{2}-\dfrac{3}{4}\right)-1$
  \task $\dfrac{2}{3}-\left[\dfrac{1}{4}-\left(\dfrac{5}{6}-1\right)\right]$
  \task $\left(1-\dfrac{1}{3}\right)-\left(\dfrac{1}{2}-\dfrac{1}{5}\right)$
  \task $\dfrac{3}{4}-\left[2-\left(\dfrac{1}{3}+\dfrac{5}{6}\right)\right]$
\end{tasks}

\ex{3} Quina fracció cal sumar a $\dfrac{3}{8}$ per a obtindre $\dfrac{5}{6}$? Quina cal restar a $\dfrac{1}{4}$ per a obtindre $-\dfrac{2}{3}$?

\ex{3} Troba l'error i corregeix-lo: \ $\dfrac{2}{3}+\dfrac{1}{5}=\dfrac{2+1}{3+5}=\dfrac{3}{8}$


% =====================================================
\seccio{8}{Multiplicació i divisió de fraccions}

\ex{1} Calcula i simplifica:\begin{tasks}[after-item-skip=6pt](3)
  \task $\dfrac{2}{3}\cdot \dfrac{5}{7}$
  \task $\dfrac{3}{2}\cdot \dfrac{2}{5}$
  \task $\dfrac{3}{7}\cdot 14$
  \task $\dfrac{4}{3}:\dfrac{5}{6}$
  \task $\dfrac{3}{5}:\dfrac{9}{10}$
  \task $6:\dfrac{3}{4}$
\end{tasks}\recorda{per a dividir, multipliquem per la \textbf{inversa}: $\dfrac{a}{b}:\dfrac{c}{d}=\dfrac{a}{b}\cdot\dfrac{d}{c}=\dfrac{a\cdot d}{b\cdot c}$.}

\ex{1} Escriu la inversa de cada nombre:\begin{tasks}(5)
  \task $\dfrac{2}{5}$
  \task $-\dfrac{7}{4}$
  \task $3$
  \task $-\dfrac{1}{6}$
  \task $1$
\end{tasks}

\ex{2} Calcula i simplifica:\begin{tasks}[after-item-skip=6pt](3)
  \task $\left(-\dfrac{8}{15}\right)\cdot \left(-\dfrac{25}{12}\right)$
  \task $\left(-\dfrac{5}{6}\right):\dfrac{10}{9}$
  \task $\dfrac{7}{10}:\left(-14\right)$
  \task $\dfrac{14}{15}\cdot \dfrac{25}{21}$
  \task $\dfrac{3}{4}\cdot \dfrac{8}{9}:\dfrac{2}{3}$
  \task $\dfrac{2}{5}:\left(\dfrac{4}{5}\cdot \dfrac{3}{8}\right)$
\end{tasks}

\ex{2} Calcula: \ a) $\dfrac{2}{3}$ de $\dfrac{9}{10}$ \quad b) la meitat de $\dfrac{4}{5}$ \quad c) el triple de $\dfrac{5}{12}$ \quad d) $\dfrac{\;\frac{3}{4}\;}{\frac{9}{2}}$

\ex{3} Per quina fracció cal multiplicar $\dfrac{4}{9}$ per a obtindre $\dfrac{2}{3}$? Entre quina fracció cal dividir $\dfrac{5}{8}$ per a obtindre 5?


% =====================================================
\seccio{9}{Potències de fraccions}

\ex{1} Calcula:\begin{tasks}[after-item-skip=6pt](4)
  \task ${\left(\dfrac{3}{4}\right)}^{3}$
  \task ${\left(-\dfrac{2}{5}\right)}^{2}$
  \task ${\left(-\dfrac{1}{3}\right)}^{3}$
  \task ${\left(\dfrac{7}{9}\right)}^{0}$
  \task ${\left(\dfrac{1}{2}\right)}^{5}$
  \task ${\left(-\dfrac{3}{2}\right)}^{4}$
  \task ${\left(-\dfrac{2}{3}\right)}^{1}$
  \task ${\left(\dfrac{5}{2}\right)}^{2}$
\end{tasks}\recorda{$\left(\dfrac{a}{b}\right)^n=\dfrac{a^n}{b^n}$. Amb base negativa, el resultat és positiu si l'exponent és parell i negatiu si és senar.}

\ex{1} Calcula i compara: \ $\left(-\dfrac{2}{3}\right)^2$, \ $-\left(\dfrac{2}{3}\right)^2$, \ $\dfrac{2^2}{3}$ \ i \ $\dfrac{-2^2}{3^2}$.

\ex{2} Escriu com una sola potència i calcula:\begin{tasks}[after-item-skip=8pt](2)
  \task ${\left(\dfrac{2}{3}\right)}^{5}\cdot {\left(\dfrac{2}{3}\right)}^{2}:{\left(\dfrac{2}{3}\right)}^{4}$
  \task ${\left(\dfrac{3}{5}\right)}^{7}:{\left(\dfrac{3}{5}\right)}^{5}$
  \task ${\left[{\left(-\dfrac{1}{2}\right)}^{3}\right]}^{2}:{\left(-\dfrac{1}{2}\right)}^{4}$
  \task ${\left(\dfrac{4}{7}\right)}^{6}:{\left(\dfrac{4}{7}\right)}^{5}$
  \task ${\left[{\left(\dfrac{1}{3}\right)}^{2}\right]}^{2}$
  \task ${\left(\dfrac{5}{4}\right)}^{3}\cdot {\left(\dfrac{5}{4}\right)}^{0}$
\end{tasks}\recorda{amb la mateixa base: $x^m\cdot x^n=x^{m+n}$, \ $x^m:x^n=x^{m-n}$, \ $(x^m)^n=x^{m\cdot n}$.}

\ex{2} Calcula de la manera més ràpida, fent servir les propietats:\begin{tasks}[after-item-skip=6pt](2)
  \task $\left(\dfrac{5}{6}\right)^4\cdot\left(\dfrac{3}{10}\right)^4\cdot4^4$
  \task $\left(\dfrac{9}{4}\right)^3:\left(\dfrac{3}{2}\right)^3$
\end{tasks}

\ex{3} Calcula fent servir les propietats de les potències (sense calcular potències amb exponents grans):\begin{tasks}[after-item-skip=10pt](2)
  \task $\left(\dfrac{3}{4}\right)^5\cdot\left(\dfrac{4}{3}\right)^3$
  \task $\left[\left(\dfrac{2}{5}\right)^3\right]^2:\left(\dfrac{2}{5}\right)^4-\left(\dfrac{1}{5}\right)^2\cdot\left(\dfrac{5}{2}\right)^0$
  \task $\left(\dfrac{6}{7}\right)^3\cdot\left(\dfrac{7}{3}\right)^3:2^2+\left(-\dfrac{1}{2}\right)^3$
  \task $\left[\left(-\dfrac{2}{3}\right)^2\right]^3\cdot\left(\dfrac{3}{2}\right)^5-\left(\dfrac{1}{2}-\dfrac{3}{4}\right)^2$
  \task $\dfrac{\left(\frac{5}{3}\right)^4\cdot\left(\frac{3}{5}\right)^2}{\left(\frac{5}{3}\right)^3}+\left(1-\dfrac{4}{5}\right)^2\cdot5^2$
  \task $\left[\left(\dfrac{1}{2}\right)^3-\left(\dfrac{1}{2}\right)^2\right]:\left(-\dfrac{1}{2}\right)^3-\left(\dfrac{4}{9}\right)^3\cdot\left(\dfrac{3}{2}\right)^6$
\end{tasks}\recorda{busca bases iguals (sumes o restes els exponents) o exponents iguals (multipliques o divideixes les bases): $\left(\dfrac{3}{4}\right)^3\cdot\left(\dfrac{4}{3}\right)^3=\left(\dfrac{3}{4}\cdot\dfrac{4}{3}\right)^3=1$.}

\ex{3} Simplifica: \ $\dfrac{\left(\frac{2}{5}\right)^3\cdot\left(\frac{2}{5}\right)^4}{\left[\left(\frac{2}{5}\right)^2\right]^3}$ \qquad i \qquad $\dfrac{\left(-\frac{3}{4}\right)^5\cdot\left(-\frac{3}{4}\right)}{\left(-\frac{3}{4}\right)^3\cdot\left(-\frac{3}{4}\right)^2}$


% =====================================================
\seccio{10}{Operacions combinades}

\ex{1} Calcula respectant la jerarquia de les operacions:\begin{tasks}[after-item-skip=8pt](3)
  \task $\dfrac{1}{2}+\dfrac{1}{3}\cdot \dfrac{3}{4}$
  \task $\left(\dfrac{1}{2}+\dfrac{1}{3}\right)\cdot \dfrac{3}{4}$
  \task $1-\dfrac{2}{3}:\dfrac{4}{9}$
  \task $\dfrac{3}{5}+\dfrac{2}{5}\cdot \dfrac{10}{3}$
  \task $\dfrac{5}{2}+\dfrac{1}{3}\cdot \dfrac{4}{5}$
  \task $\dfrac{4}{3}\cdot \left(\dfrac{2}{5}+\dfrac{1}{3}\right)$
\end{tasks}\recorda{ordre: 1) parèntesis; 2) potències; 3) multiplicacions i divisions, d'esquerra a dreta; 4) sumes i restes.}

\ex{2} Calcula:\begin{tasks}[after-item-skip=8pt](2)
  \task $\dfrac{1}{2}-{\left(\dfrac{2}{3}\right)}^{2}\cdot \dfrac{3}{4}$
  \task ${\left(\dfrac{2}{3}\right)}^{2}-\dfrac{1}{3}\cdot \left(2-\dfrac{1}{2}\right)$
  \task $\dfrac{5}{6}:\left(\dfrac{1}{2}-\dfrac{1}{3}\right)-{2}^{2}$
  \task ${\left(-\dfrac{1}{2}\right)}^{3}+\dfrac{3}{4}:\left(1-\dfrac{1}{3}\right)$
  \task ${\left(\dfrac{3}{4}-1\right)}^{2}\cdot 8+\dfrac{1}{2}$
  \task $\dfrac{21}{4}-\dfrac{2}{5}:\dfrac{4}{3}$
\end{tasks}

\ex{2} Calcula:\begin{tasks}[after-item-skip=8pt](2)
  \task $2-\left[\dfrac{1}{3}+{\left(\dfrac{1}{2}\right)}^{2}\right]:\dfrac{5}{6}$
  \task ${\left[\dfrac{3}{4}-{\left(\dfrac{1}{2}\right)}^{2}\right]}^{2}:{\left(-\dfrac{1}{2}\right)}^{3}+\dfrac{5}{3}\cdot \dfrac{6}{5}$
  \task $\left(\dfrac{1}{4}+\dfrac{2}{3}-\dfrac{1}{12}\right)\cdot 3+\left(\dfrac{3}{2}-\dfrac{5}{4}\right):\dfrac{1}{2}$
  \task $\dfrac{3}{2}:\dfrac{1}{4}\cdot \dfrac{1}{5}+\dfrac{3}{2}\cdot \dfrac{7}{5}$
\end{tasks}

\ex{3} Calcula: \ $\dfrac{\frac{2}{3}-\frac{1}{4}}{1+\frac{1}{6}}-\left(\dfrac{3}{4}-1\right)^2\cdot8$

\ex{3} Posa els parèntesis necessaris perquè cada igualtat siga certa: \ a) $\dfrac{1}{2}+\dfrac{1}{2}\cdot\dfrac{2}{3}=\dfrac{2}{3}$ \quad b) $1-\dfrac{1}{3}:\dfrac{1}{3}+1=\dfrac{1}{2}$


% =====================================================
\seccio{11}{Problemes amb fraccions}

\ex{1} Dels 28 alumnes d'una classe, $\dfrac{3}{7}$ són xics. Quants xics i quantes xiques hi ha?

\ex{1} He llegit 120 pàgines d'un llibre, que són els $\dfrac{3}{5}$ del total. Quantes pàgines té el llibre? Quantes me'n queden?

\ex{1} Un entrepà costa $3{,}60$ € i el pa en representa $\dfrac{1}{4}$; el pernil, $\dfrac{2}{3}$. Quant costa el pa? I el pernil? Quant costa la resta?

\ex{2} De tot un sopar, en Joan s'ha menjat $\dfrac{1}{7}$, l'Elena $\dfrac{2}{5}$ i la Maite $\dfrac{1}{8}$. Quina fracció s'han menjat entre tots tres? Quina fracció n'ha sobrat? Qui ha menjat més?

\ex{2} Un dipòsit de 240 L està ple. El dilluns se'n gasta $\dfrac{1}{4}$ i el dimarts, $\dfrac{2}{3}$ del que quedava. Quants litres queden? Quina fracció del dipòsit és?

\ex{2} En una votació, $\dfrac{3}{8}$ dels votants voten sí, $\dfrac{1}{3}$ voten no i la resta, 14 persones, s'abstenen. Quantes persones han votat?

\ex{2} Una ampolla plena pesa $1{,}2$ kg i, quan n'hem buidat $\dfrac{3}{4}$, pesa $0{,}48$ kg. Quant pesa l'ampolla buida?

\ex{3} La Marta gasta $\dfrac{2}{5}$ de la paga en roba i $\dfrac{1}{3}$ del que li queda en el cine. Encara li queden 24 €. Quant cobrava de paga?

\ex{3} Un ciclista fa $\dfrac{2}{7}$ d'una etapa al matí i $\dfrac{3}{5}$ del que li queda a la vesprada. Encara li falten 40 km. Quants quilòmetres té l'etapa? Quants n'ha fet al matí?

\ex{3} Un pare reparteix 180 € entre els seus tres fills: al major li dona $\dfrac{2}{5}$, al mitjà $\dfrac{5}{8}$ del que queda i al menor, la resta. Quant rep cada un? Quina fracció del total rep el menor?

\ex{3} En una pastisseria fan 150 croissants. Al matí en venen $\dfrac{3}{5}$ i a la vesprada, $\dfrac{3}{4}$ dels que queden. Els que sobren es venen a meitat de preu, a $0{,}60$ € cada un. Quants en sobren i quant es recapta per ells?

\vfill
\begin{center}\color{muted}\small\faCheckCircle[regular]\ \ Ànim! Comprova els teus resultats amb les eines interactives dels apunts de la unitat.\end{center}

\end{document}
